\begin{Solbook}{1}
 \par \noindent  \begin {enumerate} \item \( \begin {aligned}[t] A &= \dfrac {8}{\sqrt {5}-1}-2\sqrt {5}+|2-\sqrt {5}|\\ &= \dfrac {8(\sqrt {5}+1)}{(\sqrt {5}-1)(\sqrt {5}+1)}-2\sqrt {5}+\big (\sqrt {5}-2\big )\\ &= 2(\sqrt {5}+1)-2\sqrt {5}+\sqrt {5}-2 = \sqrt {5}. \end {aligned} \) \item Với $x\ge 0$, $x\ne 9$:\\ \( \begin {aligned}[t] B&=\left (\dfrac {2}{\sqrt {x}+3}+\dfrac {9-\sqrt {x}}{(\sqrt {x}-3)(\sqrt {x}+3)}\right ):\dfrac {\sqrt {x}+1}{3-\sqrt {x}}\\ &=\dfrac {2(\sqrt {x}-3)+9-\sqrt {x}}{(\sqrt {x}-3)(\sqrt {x}+3)}\cdot \dfrac {3-\sqrt {x}}{\sqrt {x}+1} =\dfrac {\sqrt {x}+3}{(\sqrt {x}-3)(\sqrt {x}+3)}\cdot \dfrac {3-\sqrt {x}}{\sqrt {x}+1}\\ &=\dfrac {-1}{\sqrt {x}+1}. \end {aligned} \) \end {enumerate}  \hfill \qedEX 
\end{Solbook}
\begin{Solbook}{2}
 \par \noindent  \begin {listEX}[2] \item \( \begin {aligned}[t] \dfrac {2x-1}{6}-\dfrac {x}{3}&\ge 4+\dfrac {5x}{3}\\ 2x-1-2x&\ge 24+10x\\ -10x&\ge 25\\ x&\le \dfrac {5}{2}. \end {aligned} \)\\ Vậy nghiệm của bất phương trình là $x\le \dfrac {5}{2}$. \item Điều kiện $x\ne 0$. \\Trừ theo vế hai phương trình, ta được: \[\begin {cases} y=3\\ \dfrac {2}{x}+2y=-2 \end {cases} \Leftrightarrow \begin {cases} y=3\\ x=-\dfrac {1}{4} \end {cases}\] Vậy hệ phương trình có nghiệm duy nhất là $\left (-\dfrac {1}{4};\,3\right )$. \end {listEX} \hfill \qedEX 
\end{Solbook}
\begin{Solbook}{3}
 \par \noindent  Gọi số quyển vở loại 80 trang và 100 trang lần lượt là $x$ và $y$ (quyển) ($x,y\in \mathbb {N}^*,x,y<550$) \par Theo đề bài, tổng số vở là 550 quyển: $x + y = 550$ \quad (1) \par Số tiền mua vở loại $80$ trang là $6000x$ (đồng). Số tiền mua vở loại $100$ trang là $7500y$ (đồng). Tổng số tiền mua hai loại vở là $3\,600\,000$ đồng: $6000x + 7500y = 3600000$ hay $4x + 5y = 2400$ \quad (2) \par Từ (1) và (2), ta có hệ phương trình: $\begin {cases} x + y = 550 \\ 4x + 5y = 2400 \end {cases}$ \par Giải hệ phương trình ta được: $x=350, y=200$ thỏa điều kiện \par Vậy số vở loại $80$ trang là $350$ quyển, số vở loại $100$ trang là $200$ quyển.  \hfill \qedEX 
\end{Solbook}
\begin{Solbook}{4}
 \par \noindent  \begin {enumerate} \item Vì $\widehat {AOB}=90^{\circ }$ nên số đo cung nhỏ $AB$ là $90^{\circ }$. Do đó \[ l=\dfrac {90}{360}\cdot 2\pi \cdot 5=\dfrac {5\pi }{2}\,\text {cm}. \] \item Diện tích hình quạt tương ứng: $S_{\text {quạt}}=\dfrac {90}{360}\pi \cdot 5^2=\dfrac {25\pi }{4}$. \\Tam giác $OAB$ vuông cân tại $O$ nên $S_{\triangle OAB}=\dfrac {1}{2}\cdot 5\cdot 5=\dfrac {25}{2}$. \\Vậy diện tích viên phân \[ S= S_{\text {quạt}}-S_{\triangle }=\dfrac {25\pi }{4}-\dfrac {25}{2}\approx 7{,}1\,\text {cm}^2. \] \end {enumerate}  \hfill \qedEX 
\end{Solbook}
\begin{Solbook}{5}
 \par \noindent  \begin {center} \begin {tikzpicture}[scale=1.2, font=\footnotesize ] \coordinate (O) at (0,0); \def \R {2} \draw (O) circle(\R ); \coordinate (A) at (-\R ,0); \coordinate (B) at (120:\R ); \coordinate (C) at (\R ,0); \path [name path=ax] (A)--++(90:2); \path [name path=bt] (B)--++({120+90}:6); \path [name intersections={of=ax and bt,by=M}]; \path [name path=btt] (B)--++(30:5); \path [name path=ct] (C)--++(90:4); \path [name intersections={of=btt and ct, by=I}]; \path [name path=bh] (B)--++($(C)-(I)$); \path [name path=ac] (A)--(C); \path [name path=ai] (A)--(I); \path [name intersections={of=bh and ac, by=K}]; \path [name intersections={of=bh and ai, by=N}]; \path [name path=ab] (A)--(B); \path [name path=om] (O)--(M); \path [name intersections={of=ab and om, by=H}]; \draw (M)--(A)--(B)--(O)--cycle (A)--(C) (K)--(M)--(I)--(C) (B)--(K)--(I) (A)--(I); \foreach \x /\g in{O/-90,A/-180/,B/90,C/0,M/180,I/0,K/-90,N/-30,H/15} {\draw [fill=white] (\x ) circle (1.5pt) +(\g :.3) node{$\x $};} \end {tikzpicture} \end {center} \begin {enumerate} \item Vì $\triangle OAB$ cân tại $O$ và $OH\perp AB$ nên $OH$ là phân giác $\Rightarrow \widehat {AOM}=\widehat {BOM}$. \\Xét $\triangle AOM$ và $\triangle BOM$ có: $OA=OB$; $OM$ là cạnh chung; $\widehat {AOM}=\widehat {BOM}$ \par Suy ra $\triangle AOM\sim \triangle BOM$ (c.g.c) suy ra $\widehat {OAM}=\widehat {OBM}$. \par Do $MA$ là tiếp tuyến tại $A$ nên $MA\perp OA$ \par Suy ra $\widehat {OAM}=90^{\circ }$ suy ra $\widehat {OBM}=90^{\circ }$ \par Suy ra $MB\perp OB$ tại $B$ suy ra $MB$ là tiếp tuyến $(O)$. \par \item $CI$ là tiếp tuyến của $(O)$ nên $CI\perp CO$. \par Lại có $BK\perp AC$ và $MA\perp AC$ nên $NK\parallel CI$, $NB\parallel MA$. \par Xét $\triangle ACI$ có $NK\parallel CI$ suy ra $\dfrac {AK}{KC}=\dfrac {AN}{NI}$ (định lí Thalès) \par Xét $\triangle IAM$ có $NB\parallel MA$ suy ra $\dfrac {MB}{BI}=\dfrac {AN}{NI}$ (định lí Thalès) \par Mà ta có $MA=MB$ (hai tiếp tuyến $MA, MB$ cắt nhau tại $M$) và $IC=IB$ (tính chất hai tiếp tuyến cắt nhau) \par Do đó $\dfrac {AK}{KC}=\dfrac {MA}{CI}$ \par Xét $\triangle AKM \backsim \triangle CKI$ có $\dfrac {AK}{KC}=\dfrac {MA}{CI}$; $\widehat {KAM}=\widehat {KCI}(=90^\circ )$ \par Suy ra $\triangle AKM\sim \triangle CKI \text { (c.g.c)}$ Suy ra $\widehat {AKM}=\widehat {CKI}$ hay $\widehat {BKM}=\widehat {BKI}$ \par Suy ra $KB$ là tia phân giác của $\widehat {MKI}$. \end {enumerate}  \hfill \qedEX 
\end{Solbook}
\begin{Solbook}{6}
 \par \noindent  \begin {enumerate} \item ĐKXĐ: $x\geq 1$ \\ \( x^{2}+22x-7=8(x+3)\sqrt {x-1}\\ x^{2}+22x-7-8(x+3)\sqrt {x-1}=0\\ x^{2}+6x+9-8(x+3)\sqrt {x-1}+16(x-1)=0\\ (x+3-4\sqrt {x-1})^{2}=0\\ x+3=4\sqrt {x-1}\\ x^{2}+6x+9=16x-16 \quad (x \geq 1)\\ x^{2}-10x+25=0\\ x=5 \text { (thỏa mãn ĐKXĐ).} \)\\ Vậy nghiệm của phương trình là $x=5$. \item Với $a, b$ là các số dương ta có $(\sqrt {a}-\sqrt {b})^{2} \geq 0$ suy ra $\sqrt {a} \cdot \sqrt {b} \leq \dfrac {1}{2}(a+b)$ $(*)$ \\ Dấu "=" xảy ra khi $a=b$. \par Theo đề bài $x, y, z$ là số dương và $x+y+z=2025 x y z$ nên $\dfrac {2025 x y z}{x+y+z}=1$ \\ Do đó $\sqrt {1+2025 x^{2}}=\sqrt {\dfrac {2025 x y z}{x+y+z}+2025 x^{2}} =\sqrt {2025x \dfrac {yz+x^{2}+xy+xz}{2025xyz}}=\sqrt {\dfrac {(x+y)(y+z)}{yz}}$ Suy ra \[ \dfrac {1}{\sqrt {1+2025 x^{2}}}=\sqrt {\dfrac {y z}{(x+y)(y+z)}}=\sqrt {\dfrac {y}{x+y}} \cdot \sqrt {\dfrac {z}{x+z}} \leq \dfrac {1}{2}\left (\dfrac {y}{x+y}+\dfrac {z}{x+z}\right ) \text { (Áp dụng BĐT $(*)$)} \] Dấu "=" xảy ra khi $\dfrac {y}{x+y}=\dfrac {z}{x+z}$ hay $y=z$. \par Tương tự: \[ \dfrac {1}{\sqrt {1+2025 y^{2}}}=\sqrt {\dfrac {x z}{(y+x)(y+z)}}=\sqrt {\dfrac {x}{y+x}} \cdot \sqrt {\dfrac {z}{y+z}} \leq \dfrac {1}{2}\left (\dfrac {x}{x+y}+\dfrac {z}{y+z}\right ) \] Dấu "=" xảy ra khi $x=z$ \[ \dfrac {1}{\sqrt {1+2025 z^{2}}}=\sqrt {\dfrac {x y}{(z+x)(z+y)}}=\sqrt {\dfrac {x}{x+z}} \cdot \sqrt {\dfrac {y}{y+z}} \leq \dfrac {1}{2}\left (\dfrac {x}{x+z}+\dfrac {y}{y+z}\right ) \] Dấu "=" xảy ra khi $x=y$ \par Nên \[ P \leq \dfrac {1}{2}\left (\dfrac {y}{x+y}+\dfrac {z}{x+z}+\dfrac {x}{x+y}+\dfrac {z}{y+z}+\dfrac {x}{x+z}+\dfrac {y}{y+z}\right )=\dfrac {3}{2} \] Dấu "=" xảy ra khi $x=y=z=\dfrac {\sqrt {3}}{45}$. \\ Vậy giá trị lớn nhất của $P$ là $\dfrac {3}{2}$, đạt được khi $x=y=z=\dfrac {\sqrt {3}}{45}$. \end {enumerate} \hfill \qedEX 
\end{Solbook}
